\documentclass[preprint,12pt,authoryear]{elsarticle}

\usepackage[utf8]{inputenc}
\usepackage{lmodern}
\usepackage{amssymb}
\usepackage{amsmath}
\usepackage[hidelinks]{hyperref}  % hidelinks = no visible coloured boxes around \ref / \cite
\usepackage{xurl}  % allow URL line-breaking anywhere (must load after hyperref)
\usepackage{booktabs}
\usepackage{graphicx}
\usepackage{float}
\usepackage{multirow}
\usepackage{enumitem}
\usepackage{ragged2e}  % \justifying for table-note paragraphs inside centred floats
\setlength{\emergencystretch}{6em}

\makeatletter
\pdfstringdefDisableCommands{%
  \def\corref#1{}%
  \def\@corref#1{}%
  \def\cnotenum{}%
}
\makeatother

\journal{Finance Research Letters}

\newcommand{\SPY}{SPY}
\newcommand{\CL}{CL}
\newcommand{\USDJPY}{USD/JPY}
\newcommand{\GLD}{GLD}

% Supplement-only: left-align ONLY the title; authors/affiliations stay
% centered (elsarticle preprint default).  Also drop the two \hrule
% lines that bracket the (empty) abstract slot.
\makeatletter
\long\def\pprintMaketitle{\clearpage
  \iflongmktitle\if@twocolumn\let\columnwidth=\textwidth\fi\fi
  \resetTitleCounters
  \def\baselinestretch{1}%
  \printFirstPageNotes
  \thispagestyle{pprintTitle}%
  \def\baselinestretch{1}%
  \begin{flushleft}%
    \Large\@title\par
  \end{flushleft}%
  \vskip18pt
  \ifx\@elsarticlenewpageafter\newpage@after@title
    \newpage
  \fi
  \begin{center}%
    \ifdoubleblind
      \vspace*{2pc}
    \else
      \normalsize\elsauthors\par\vskip10pt
      \footnotesize\itshape\elsaddress\par\vskip18pt
    \fi
    \ifx\@elsarticlenewpageafter\newpage@after@author
      \newpage
    \fi
    \ifvoid\absbox\else\unvbox\absbox\par\vskip10pt\fi
    \ifvoid\keybox\else\unvbox\keybox\par\vskip10pt\fi
    \ifx\@elsarticlenewpageafter\newpage@after@abstract
      \newpage
    \fi
  \end{center}%
  \gdef\thefootnote{\arabic{footnote}}%
}
\makeatother

\begin{document}

\begin{frontmatter}
\sloppy

\title{{\large Supplement}\\[8pt] Regime Labels Are Not Resolution-Invariant: Evidence from the 2026 US--Iran Escalation}

\author[1]{Kai Zheng}
\ead{kaizhengnz@gmail.com}
\ead[orcid]{0009-0004-6661-2023}
\author[2]{Rand Low}
\ead{rlow@bond.edu.au}
\ead[orcid]{0000-0002-0290-7840}
\author[1]{Ruili Wang\corref{cor1}}
\ead{ruili.wang@massey.ac.nz}
\ead[orcid]{0000-0003-2899-9816}

\cortext[cor1]{Corresponding author: Ruili Wang}

\affiliation[1]{organization={School of Mathematical and Computational Sciences\\Massey University},
            city={Auckland},
            country={New Zealand}}

\affiliation[2]{organization={Bond Business School\\Bond University},
            city={Gold Coast},
            state={QLD},
            country={Australia}}

\end{frontmatter}
\sloppy  % relax line-breaking globally to absorb dense citation/math runs without overfull-hbox

% S-prefixed numbering for sections, tables, figures, equations
\renewcommand{\thesection}{S\arabic{section}}
\renewcommand{\thetable}{S\arabic{table}}
\renewcommand{\thefigure}{S\arabic{figure}}
\renewcommand{\theequation}{S\arabic{equation}}

The following tables and analyses constitute the Supplement accompanying the main manuscript.

\section{Data and instrument-choice justifications}\label{app:data_choices}

\paragraph{Choice of base resolution}
We anchor on the 5-minute frequency for two reasons. First, below 5m bid--ask bounce and stale-quote effects increasingly dominate realised-volatility estimators on the assets in our universe \citep{AitSahaliaMykland2005,Andersen1999RealizedVolatility}, motivating 5m as the floor; above-floor identification of two-component GMMs is then asset-specific (\S{}3.1, Table~\ref{tab:gmm_diag}). Second, the diagnostic question of this paper is whether \emph{labels} agree across the resolutions a desk actually uses for limit-setting, monitoring, and reporting (typically intraday-to-daily); going below 5m would generate cross-frequency pairs that no production risk system consumes. IB \texttt{CONTFUT} 1-minute history at our required panel length is also rate-limited; this is a binding but secondary operational constraint rather than an independent methodological argument.

\paragraph{Choice of equity instrument}
We use SPDR S\&P 500 ETF (SPY) rather than the S\&P 500 spot index because the index is non-tradable, has no executable bid/ask or volume, and is itself a stale-quote-prone synthetic of constituent prints \citep{HasbrouckOneSecurity1995}. SPY is the most liquid US equity ETF and the instrument on which limits, hedges, and SR~26-2 model documentation are actually written for US equity desks; the 5m SPY bar is therefore the operationally correct unit of analysis for a regime-classifier diagnostic.

\paragraph{Choice of FX pair}
The umbrella rationale for selecting USD/JPY is informativeness about a distinct microstructure, not representativeness of the FX market. Two features make USD/JPY the informative FX probe here: (i)~JPY is the canonical safe-haven leg of geopolitical risk transmission \citep{CaldaraIacoviello2022}, so the 2022 and 2026 episodes load directly onto the pair (EUR/USD reacts only through second-order channels: oil pass-through, ECB policy reaction); (ii)~USD/JPY's tri-session microstructure (Tokyo / London / New York) generates session-boundary heterogeneity in volatility that EUR/USD's bi-session structure does not, providing a more demanding cross-frequency stress test for the regime classifier. The pair is therefore informativeness-selected, not representativeness-selected: results should be read as one geopolitically- and tri-session-loaded FX probe, not an FX-market average. 
\paragraph{Continuous-futures construction}
The continuous CL contract is sourced via Interactive Brokers' \texttt{CONTFUT} identifier, which delivers a continuous front-month series adjusted at rollover via IB's proprietary multiplicative (ratio) method to maintain price continuity across contract changes; we do not have access to the unadjusted contract-by-contract series under this identifier. Log-returns on a ratio-adjusted continuous series are unbiased relative to within-contract returns at non-roll dates, preserving the validity of vol-based GMM classification away from the excluded roll window. The roll window is approximated as calendar days 18--22 of each month, corresponding to the typical WTI NYMEX roll period (CME WTI final settlement falls on the third business day before the 25th of the prior calendar month, which lies between the 18th and 22nd in calendar days); bars in this window are excluded from the rolling-volatility estimate to mitigate artificial spikes. Empirical impact is small: roll-week ARI is 0.388 vs.\ 0.392 in non-roll weeks (Table~\ref{tab:cl_roll}), computed on the GMM-with-percentile-fallback labels described in \S{}2.2 (CL 2026 5m/15m fits are degenerate, \S{}3.1, Table~\ref{tab:gmm_diag}). The comparison is therefore informative about roll-window mechanical effects rather than about GMM-classifier identification per se; the directional implication (roll-window inclusion does not materially change the headline ARI) rules out roll mechanics as a confound.
\paragraph{Inclusion of GLD in the 2022 universe}
The 2022 universe drops \CL{}: IB \texttt{CONTFUT}'s pre-2024 5m history is forward-fill placeholder data (downloaded bars show $\sim 99.6\%$ identical-OHLC across the Dec 2021--May 2022 window, with mean Volume $< 1$), unusable for cross-frequency regime analysis. SPDR Gold Trust (GLD) is added as a precious-metal commodity proxy responding to the invasion-driven safe-haven flow (gold rallied over the Dec 2021 -- early-March 2022 window, peaking 8 March 2022). The 2022 OOS sample is therefore three-asset (SPY, USD/JPY, GLD); CL replication for 2022 is left for future work when extended historical 5m data become available from a non-IB source, and QQQ 2022 5m bars were not pulled in this revision round. The 2026 universe carries all five assets (SPY, QQQ, USD/JPY, CL, GLD), giving a five-by-four cross-frequency panel against which every numerical claim in \S{}3.1--\S{}3.4 is computed; the 2022 sample is a partial replication on three assets, treated as a same-asset hold-out for SPY and USD/JPY and a within-class commodity substitution for the third probe.

\paragraph{Data treatments not applied}
For transparency, the following standard pre-processing steps are \emph{not} applied:
(i)~outlier removal beyond the frequency-aware Winsorize discussed in \S{}2.2 ($\pm 3\%/5\%/10\%/20\%$ at 5m/15m/1h/1d; the schedule replaces the earlier flat $\pm 3\%$ cap that produced a frequency-asymmetric clip rate at 1d);
(ii)~holiday-calendar adjustment for cross-asset trading-day mismatch (USD/JPY trades on US holidays where SPY does not; daily bars use the 16:00 NY close on each asset's available days, and ARI is computed on the intersection of valid days);
(iii)~explicit daylight-saving alignment at session boundaries (calendar-time bins are computed in NY local time, so the Tokyo session start drifts by one hour at DST transitions; the 15m$\to$1h single-step drop is intraday-only and not sensitive to a single-hour calendar drift);
(iv)~ETF distribution adjustment for SPY and GLD (quarterly distributions are well within the per-resolution Winsorize thresholds and are not separately handled).

\subsection{Positioning relative to the multiscale-regime literature}\label{app:lit_comparison}
The list below contrasts this paper with the three closest cousins to our object of study, namely multiscale-regime studies that operate on volatility / jump / cascade structures across time scales. The intent is to make explicit that we operate on a different unit of analysis (post-classifier discrete labels), pursue a different inferential target (operational model risk inherited at the model-consumption layer), and propose a different deliverable (a model-validation diagnostic rather than a structural decomposition).

\begin{description}\setlength\itemsep{4pt}
\item[\citet{GencaySelcukGradojevicWhitcher2010}.] Wavelet decomposition with asymmetric volatility regimes; documents asymmetric high$\to$low information flow on continuous wavelet coefficients across investment horizons. \emph{Difference:} jointly estimates all scales, smoothing the post-classifier label disagreement that the present paper measures by design.

\item[\citet{ErdemliogluGradojevic2021}.] HMM regime-switching with realised jumps; documents horizon-dependent jump-tail regimes via a single HMM shared across horizons. \emph{Difference:} the present paper fits an independent classifier at each resolution, so cross-resolution disagreement is exposed rather than absorbed into one shared transition kernel.

\item[\citet{GradojevicTsiakas2021}.] Wavelet hidden Markov tree on cryptocurrency volatility cascades; documents asymmetric short$\to$long-horizon volatility cascade through continuous cascade transition probabilities. \emph{Difference:} the object is continuous structural cascade; the present paper's object is the binary post-classifier label sequence that downstream risk systems consume.

\item[\textbf{This paper.}] Independent two-component GMM (HMM in robustness) fit at four resolutions on the same price feed; cross-frequency label agreement is measured via ARI/AMI on the discrete crisis/calm sequence. The headline finding is a uniform within-intraday excess of $+0.03$ to $+0.24$ in mean off-diagonal ARI over a canonical-pipeline calibrated $K{=}2$ MS-Gaussian baseline; the deliverable is an operational model-validation diagnostic (the \emph{resolution envelope}) mapped to SR~26-2 governance.
\end{description}

\noindent The differences are substantive, not merely methodological: by fitting a two-state classifier independently at each resolution and measuring ARI on the post-classifier label sequence, the present paper recovers the operational artefact that a desk inherits when it calibrates at one frequency and reports at another, an object that joint-estimation wavelet/HMM frameworks elide by construction. The three cited studies provide complementary structural descriptions of the underlying volatility process; the present paper translates that structural property into a model-risk diagnostic with an SR~26-2 governance hook.

% source: outputs/{SPY,CL,USDJPY,GLD}_cross_freq_ari.csv  (4x4 matrix, index_col=0)
\begin{table}[!htbp]
\centering
\small
\setlength{\tabcolsep}{4pt}
\caption{Cross-frequency ARI matrices by asset (full sample).}
\label{tab:ari_matrix_combined}
\begin{tabular}{llcccc}
\toprule
Asset & & 5m & 15m & 1h & 1d \\
\midrule
\multirow{4}{*}{\SPY{}}
 & 5m  & 1.00 & 0.539 & 0.132 & 0.068 \\
 & 15m &      & 1.00 & 0.138 & 0.051 \\
 & 1h  &      &      & 1.00 & 0.171 \\
 & 1d  &      &      &      & 1.00 \\
\midrule
\multirow{4}{*}{QQQ}
 & 5m  & 1.00 & 0.536 & 0.123 & 0.024 \\
 & 15m &      & 1.00 & 0.133 & 0.042 \\
 & 1h  &      &      & 1.00 & 0.086 \\
 & 1d  &      &      &      & 1.00 \\
\midrule
\multirow{4}{*}{\CL{}}
 & 5m  & 1.00 & 0.694 & 0.347 & 0.315 \\
 & 15m &      & 1.00 & 0.308 & 0.306 \\
 & 1h  &      &      & 1.00 & 0.393 \\
 & 1d  &      &      &      & 1.00 \\
\midrule
\multirow{4}{*}{\USDJPY{}}
 & 5m  & 1.00 & 0.409 & 0.121 & $-$0.017 \\
 & 15m &      & 1.00 & 0.183 & $-$0.025 \\
 & 1h  &      &      & 1.00 & $-$0.043 \\
 & 1d  &      &      &      & 1.00 \\
\midrule
\multirow{4}{*}{\GLD{}}
 & 5m  & 1.00 & 0.712 & 0.217 & $-$0.056 \\
 & 15m &      & 1.00 & 0.188 & $-$0.053 \\
 & 1h  &      &      & 1.00 & $-$0.079 \\
 & 1d  &      &      &      & 1.00 \\
\bottomrule
\end{tabular}
\par\medskip
{\footnotesize\justifying
\textit{Notes:} GMM on log(volatility). 5m--15m pairs agree moderately; any pair crossing the 15m$\to$1h boundary drops to near zero.
\par}
\end{table}

% source: outputs/{ASSET}_cross_freq_ari.csv         (full sample)
%         outputs/{ASSET}_event_cross_freq_ari.csv   (2026-02-28 to 04-15)
%         outputs/{ASSET}_calm_cross_freq_ari.csv    (2026-01-01 to 24)
\begin{table}[!htbp]
\centering
\small
\setlength{\tabcolsep}{4pt}
\caption{Per-asset cross-frequency ARI by sub-window: full sample / event window (2026-02-28 to 04-15) / calm window (2026-01-01 to 2026-01-24). The triple-cell breakdown isolates the within-period dissonance from the cross-period mean-shift coherence: CL's full-sample 1h--1d 0.393 collapses to event 0.00 / calm 0.00, identifying the headline coherence as event-driven volatility-level shift detection rather than within-period label agreement.}
\label{tab:cross_freq_subwindow}
\begin{tabular}{llcccc}
\toprule
Asset & Window & 5m--15m & 15m--1h & 1h--1d & All-6 mean \\
\midrule
\multirow{3}{*}{SPY}
  & Full   & 0.539 & 0.138            & 0.171    & 0.183 \\
  & Event  & 0.567 & 0.132            & 0.298    & 0.202 \\
  & Calm   & 0.434 & 0.015            & 0.041    & 0.097 \\[3pt]
\multirow{3}{*}{QQQ}
  & Full   & 0.536 & 0.133            & 0.086    & 0.157 \\
  & Event  & 0.521 & 0.125            & 0.056    & 0.142 \\
  & Calm   & 0.427 & 0.034            & $-$0.010 & 0.082 \\[3pt]
\multirow{3}{*}{USD/JPY}
  & Full   & 0.409 & 0.183$^\dagger$  & $-$0.043 & 0.105 \\
  & Event  & 0.332 & 0.109$^\dagger$  & $-$0.082 & 0.052 \\
  & Calm   & 0.423 & 0.079$^\dagger$  & $-$0.027 & 0.079 \\[3pt]
\multirow{3}{*}{CL}
  & Full   & 0.694$^\dagger$ & 0.308$^\dagger$  & 0.393    & 0.394 \\
  & Event  & 0.464$^\dagger$ & 0.002$^\dagger$  & $-$0.010 & 0.079 \\
  & Calm   & 0.451$^\dagger$ & 0.240$^\dagger$  & 0.000    & 0.151 \\[3pt]
\multirow{3}{*}{GLD}
  & Full   & 0.712 & 0.188            & $-$0.079 & 0.155 \\
  & Event  & 0.694 & 0.125            & $-$0.018 & 0.141 \\
  & Calm   & 0.804 & 0.059            & $-$0.052 & 0.125 \\
\bottomrule
\end{tabular}
\par\medskip
{\footnotesize\justifying
\textit{Notes:} Event window = 2026-02-28 to 2026-04-15 (Iran joint-strike to ceasefire). Calm window = 2026-01-01 to 2026-01-24 (pre-escalation). The 5m--15m / 15m--1h pairs are essentially window-invariant for SPY, QQQ, USD/JPY, GLD: dissonance at those resolution gaps is structural, not stress-driven. CL is qualitatively different at 1h--1d: full-sample 0.393 vs.\ event 0.00 vs.\ calm 0.00, so the headline coherence comes from cross-period mean-shift detection (\S{}3.1, ``CL is event coherence, not Layer 2 invariance''), not from within-period day-by-day agreement. SPY 1h--1d shows mild stress amplification (calm 0.04 $\to$ event 0.30, full 0.17): SPY 1h and 1d agree more on which days are stress-period crisis days than on which are calm-period crisis days, but the event-only reading is still well below the 5m--15m baseline. $\dagger$ Pairs touching a percentile-fallback endpoint (USD/JPY 15m, CL 5m / 15m); see \S{}2.2, Suppl Table~\ref{tab:gmm_diag}.
\par}
\end{table}

% source: qualitative summary; numeric inputs from outputs/{ASSET}_daily_summary.csv (date=2026-03-06)
\begin{table}[!htbp]
\centering
\small
\setlength{\tabcolsep}{4pt}
\caption{6 March 2026: qualitative interpretation of cross-resolution conflict.}
\label{tab:mar06_conflict}
\begin{tabular}{lp{3.5cm}p{7.5cm}}
\toprule
Asset & Pattern & Interpretation \\
\midrule
SPY & all resolutions crisis & Consistent classification on this day; the inversion mechanism does not fire here, but \SPY{} disagrees on 32\% of full-sample days (Table~\ref{tab:var_uplift}). \\
CL & high-vol intraday, 1d calm & 5m and 15m bars show high-vol via percentile-fallback labels (GMM degenerate, \S{}2.2); 1h is GMM-identified at 100\% crisis; 1d is calm, consistent with the news-event narrative of intraday supply-shock spikes that mean-revert within the session, though the causal interpretation rests on the news event rather than a within-paper return-path test. \\
USD/JPY & partial-crisis intraday (27--34\%), 1d crisis & Intensification rather than inversion: a partial-crisis intraday slice lifts to outright crisis at the daily horizon as overnight repricing across Tokyo/London/NY sessions is absorbed into the close-to-close bar but averaged away inside any single hourly slice. \\
GLD & partial-crisis intraday (23--24\%), 1d crisis & Same intensification mechanism as \USDJPY{}; safe-haven inflows arrive at session boundaries. \\
\bottomrule
\end{tabular}
\par\medskip
{\footnotesize\justifying
\textit{Notes:} Crisis shares from Table~2; GMM on log(volatility).
\par}
\end{table}

% source: outputs/{ASSET}_cross_freq_ari.csv             (Baseline)
%         outputs/{ASSET}_tod_adjusted_cross_freq_ari.csv (TOD-adjusted)
%         outputs/{ASSET}_calendar_window_cross_freq_ari.csv (Calendar-aligned, intraday-only mean)
\begin{table}[!htbp]
\centering
\small
\setlength{\tabcolsep}{5pt}
\caption{Decomposition of resolution dissonance: mean off-diagonal ARI under baseline, TOD-adjusted, and calendar-time-aligned specifications.}
\label{tab:decomposition}
\begin{tabular}{lccc}
\toprule
Asset & Baseline & TOD-adjusted & Calendar-aligned (intraday) \\
\midrule
SPY & 0.270 & 0.266 & 0.369 \\
QQQ & 0.264 & 0.183 & 0.367 \\
CL & 0.450 & 0.462 & 0.670 \\
USD/JPY & 0.237 & 0.239 & 0.144 \\
GLD & 0.372 & 0.417 & 0.493 \\
\bottomrule
\end{tabular}
\par\medskip
{\footnotesize\justifying
\textit{Notes:} All three columns report mean off-diagonal ARI over intraday pairs only (5m--15m, 5m--1h, 15m--1h); 1d is excluded because the 6-hour calendar-aligned window yields degenerate estimates at daily frequency. Baseline values therefore differ from the full 4$\times$4 means in Table~1, which include daily pairs. Baseline = frequency-specific rolling windows (2h for 5m/15m, 24 bars for 1h). TOD-adjusted = volatility normalised by hourly median before GMM fitting. Calendar-aligned = fixed 6-hour rolling window for all intraday frequencies. \emph{Look-ahead caveat (TOD-adjusted column only):} the per-hour normalisation uses the full-sample hourly median; the TOD-adjusted column is therefore an in-sample diagnostic and should not be read as a no-look-ahead robustness check.
\par}
\end{table}

% source: outputs/pipeline_summary.csv  (overall_mean_ari)
%         outputs/hypothesis_tests.csv  (perm p-values, null 95% interval)
\begin{table}[!htbp]
\centering
\small
\setlength{\tabcolsep}{6pt}
\caption{Permutation test for full-sample mean off-diagonal ARI (overall).}
\label{tab:perm_overall}
\begin{tabular}{lccc}
\toprule
Asset & Mean off-diag.\ ARI & Permutation $p$-value & Null 95\% interval \\
\midrule
SPY & 0.183 & $<$0.002 & [$-$0.001, 0.000] \\
QQQ & 0.157 & $<$0.002 & [$-$0.001, 0.001] \\
CL & 0.394 & $<$0.002 & [$-$0.002, 0.002] \\
USD/JPY & 0.105 & $<$0.002 & [$-$0.001, 0.002] \\
GLD & 0.155 & $<$0.002 & [$-$0.003, 0.003] \\
\bottomrule
\end{tabular}
\par\medskip
{\footnotesize\justifying
\textit{Notes:} 500 permutations per asset. The observed statistic exceeded all permutations in every case, so $p < 0.002$ (= $1/501$ with add-one smoothing). Higher-resolution $p$-values would require $\ge$5{,}000 permutations. Null: independently permute each resolution's label sequence, preserving marginal crisis shares (\emph{weak} independently-shuffled-labels null; this is a sanity check, not primary evidence). 5-asset 2026 mean = 0.199.
\par}
\end{table}

% source: outputs/robustness_summary.csv  (baseline ARI by asset)
%         outputs/robustness_ranges.csv   (sweep min-max over K x window_mult)
\begin{table}[!htbp]
\centering
\small
\setlength{\tabcolsep}{4pt}
\caption{Robustness to alternative $K$ and rolling-volatility windows.}
\label{tab:robustness}
\resizebox{\textwidth}{!}{%
\begin{tabular}{lcccc}
\toprule
Asset & Baseline full-sample ARI & Sweep range & Baseline latest 7-day ARI & Sweep range \\
\midrule
SPY & 0.183 & 0.082--0.204 & 0.087 & 0.046--0.240 \\
QQQ & 0.157 & 0.073--0.157 & 0.100 & 0.064--0.182 \\
CL & 0.394 & 0.153--0.541 & 0.100 & 0.069--0.154 \\
USD/JPY & 0.105 & 0.037--0.252 & 0.276 & 0.077--0.415 \\
GLD & 0.155 & 0.044--0.350 & 0.257 & $-$0.019--0.318 \\
\bottomrule
\end{tabular}%
}
\par\medskip
{\footnotesize\justifying
\textit{Notes:} Baseline = $K{=}2$ with 1.0$\times$ volatility windows (2h for 5m/15m, 24 bars for 1h, 5 bars for 1d). Sweep varies $K \in \{2,3\}$ and the window multiplier in \{0.5, 1.0, 2.0\}. \CL{}'s sweep upper bound of 0.541 is reached at $K{=}3$, $2.0\times$ window: the third state captures supply-shock days that the binary classifier shuffles between calm and crisis, raising same-direction agreement at the supply-shock boundary; the underlying 15m$\to$1h single-step drop is preserved at every cell of the sweep (Table~\ref{tab:k3_baseline}).
\par}
\end{table}

\paragraph{Parameter extremes}
Table~\ref{tab:robustness_extremes} reports the specifications producing the lowest and highest full-sample ARI.

% source: outputs/robustness_ranges.csv  (min/max ARI with corresponding K, window mult)
\begin{table}[!htbp]
\centering
\small
\setlength{\tabcolsep}{4pt}
\caption{Parameter combinations delivering the minimum and maximum full-sample ARI (2026 panel).}
\label{tab:robustness_extremes}
\resizebox{\textwidth}{!}{%
\begin{tabular}{lcccccc}
\toprule
Asset & Min ARI & $K$ & Window mult. & Max ARI & $K$ & Window mult. \\
\midrule
SPY & 0.082 & 3 & 2.0$\times$ & 0.204 & 2 & 2.0$\times$ \\
QQQ & 0.073 & 3 & 2.0$\times$ & 0.157 & 2 & 1.0$\times$ \\
CL & 0.153 & 2 & 0.5$\times$ & 0.541 & 3 & 2.0$\times$ \\
USD/JPY & 0.037 & 3 & 0.5$\times$ & 0.252 & 3 & 2.0$\times$ \\
GLD & 0.044 & 2 & 0.5$\times$ & 0.350 & 3 & 2.0$\times$ \\
\bottomrule
\end{tabular}%
}
\medskip
\footnotesize
\end{table}

\paragraph{K=3 baseline-window mean ARI (tri-state classifier)}
Table~\ref{tab:k3_baseline} reports the per-asset mean off-diagonal ARI under $K{=}3$ at the baseline window multiplier (1.0$\times$), separating 2026 and 2022 OOS. The hierarchical pattern is preserved on every cell: 5m--15m moderate, 5m--1h and 15m--1h $\le 0.20$, intraday-to-daily near zero. \CL{} 2026 sits at 0.390 because the third state captures the supply-shock days that the binary classifier flips between calm and crisis, raising same-resolution agreement; the structural-break across 15m$\to$1h is qualitatively unchanged (Table~\ref{tab:robustness}).

% source: outputs/robustness_summary.csv      (filter K=3, window_mult=1.0; 2026)
%         outputs_2022/robustness_summary.csv (filter K=3, window_mult=1.0; 2022 OOS)
\begin{table}[H]
\centering
\small
\setlength{\tabcolsep}{6pt}
\caption{K=3 baseline-window (1.0$\times$) mean off-diag.\ ARI by asset and episode.}
\label{tab:k3_baseline}
\begin{tabular}{lccc}
\toprule
Asset & 2026 & 2022 OOS & Combined range \\
\midrule
SPY & 0.160 & 0.249 & 0.16--0.25 \\
QQQ & 0.140 & --- & --- \\
USD/JPY & 0.082 & 0.241 & 0.08--0.24 \\
CL & 0.390 & --- & --- \\
GLD & 0.229 & 0.078 & 0.08--0.23 \\
\bottomrule
\end{tabular}
\par\medskip
{\footnotesize\justifying
\textit{Notes:} GMM with $K{=}3$ on log(volatility); the high-mean state is treated as crisis, the lowest as calm, and the middle state as a calm-adjacent intermediate (collapsing middle-vs-low produces a binary label compatible with $K{=}2$). Window multiplier fixed at 1.0$\times$ (2h for 5m/15m, 24 bars for 1h, 5 bars for 1d). 5-asset 2026 mean = 0.200; 3-asset 2022 mean = 0.189. Both empirical means exceed the corresponding $K{=}2$ \emph{empirical} mean (Table~\ref{tab:robustness}, full-sample) because the third state absorbs supply-shock days that the binary classifier shuffles. The canonical-pipeline $K{=}3$ calibrated baseline is 0.113 with simulation IQR $[0.097, 0.125]$ at the ML-fit persistence (\S\ref{app:k3_baseline}); the empirical $K{=}3$ five-asset 2026 mean exceeds the calibrated upper bound. \emph{Per-asset disaggregation:} CL 2026 (0.390) is the largest excess (third state absorbs CL's bursty supply-shock days); SPY 2022 (0.249), USD/JPY 2022 (0.240) and GLD 2026 (0.229) similarly exceed the calibrated IQR upper bound; SPY 2026 (0.160) and QQQ 2026 (0.140) are above the upper bound by 0.015--0.035; USD/JPY 2026 (0.082) and GLD 2022 (0.078) lie below the calibrated IQR. The within-cell hierarchical pattern (5m--15m high, 15m$\to$1h low, 1h--1d near zero) is unchanged from $K{=}2$.
\par}
\end{table}

\paragraph{GMM fit diagnostics}
Table~\ref{tab:gmm_diag} reports BIC, component means, separation, and overlap for each asset$\times$frequency combination.

% source: outputs/{SPY,CL,USDJPY,GLD}_gmm_diagnostics.csv  (one row per asset x freq)
\begin{table}[!htbp]
\centering
\small
\setlength{\tabcolsep}{4pt}
\caption{GMM fit diagnostics by asset and resolution (2026 primary sample).}
\label{tab:gmm_diag}
\resizebox{\textwidth}{!}{%
\begin{tabular}{llrrrrrr}
\toprule
Asset & Freq & $n$ & BIC & $\mu_{\text{calm}}$ & $\mu_{\text{crisis}}$ & Separation & Overlap \\
\midrule
SPY & 5m & 23,688 & 50,130.1 & $-$8.10 & $-$7.31 & 1.39 & 0.487 \\
SPY & 15m & 8,020 & 18,564.7 & $-$7.71 & $-$6.86 & 1.32 & 0.511 \\
SPY & 1h & 2,098 & 2,790.0 & $-$6.54 & $-$6.06 & 1.18 & 0.558 \\
SPY & 1d & 124 & 212.9 & $-$5.30 & $-$4.65 & 1.59 & 0.426 \\
\midrule
QQQ & 5m & 23,688 & 50,175.6 & $-$7.84 & $-$7.01 & 1.48 & 0.459 \\
QQQ & 15m & 8,020 & 18,462.6 & $-$7.44 & $-$6.56 & 1.40 & 0.486 \\
QQQ & 1h & 2,098 & 2,473.8 & $-$6.21 & $-$5.74 & 1.27 & 0.523 \\
QQQ & 1d & 124 & 219.7 & $-$5.03 & $-$4.49 & 1.14 & 0.582 \\
\midrule
CL & 5m & 34,896 & 72,824.5 & $-$6.70 & $-$27.63 & \textbf{degen.} & $\approx$0 \\
CL & 15m & 11,759 & 27,574.9 & $-$6.22 & $-$27.63 & \textbf{degen.} & $\approx$0 \\
CL & 1h & 3,037 & 5,108.5 & $-$5.74 & $-$4.73 & 3.05 & 0.130 \\
CL & 1d & 155 & 398.3 & $-$4.21 & $-$3.08 & 1.80 & 0.378 \\
\midrule
USD/JPY & 5m & 36,393 & 49,987.1 & $-$8.50 & $-$8.33 & 0.36 & 0.850 \\
USD/JPY & 15m & 12,259 & 19,863.2 & $-$7.95 & $-$27.63 & \textbf{degen.} & $\approx$0 \\
USD/JPY & 1h & 3,067 & 2,223.1 & $-$7.20 & $-$6.84 & 1.13 & 0.584 \\
USD/JPY & 1d & 156 & 245.1 & $-$5.99 & $-$5.49 & 1.16 & 0.559 \\
\midrule
GLD & 5m & 23,688 & 47,123.4 & $-$7.19 & $-$6.52 & 1.08 & 0.591 \\
GLD & 15m & 8,020 & 17,674.1 & $-$6.74 & $-$6.27 & 0.67 & 0.733 \\
GLD & 1h & 2,098 & 2,923.6 & $-$5.86 & $-$5.24 & 1.49 & 0.453 \\
GLD & 1d & 124 & 278.3 & $-$4.85 & $-$4.06 & 1.29 & 0.536 \\
\bottomrule
\end{tabular}%
}
\par\medskip
{\footnotesize\justifying
\textit{Notes:} GMM with $K{=}2$ on log(volatility). $\mu$ = component means in log-vol space. For non-degenerate cells the lower-mean component is shown in the $\mu_{\text{calm}}$ column (lower log-volatility = calm, the standard finance convention). For \textbf{degen.} cells the dominant-mass component (always at $\mu \approx -6$ to $-8$) is shown in the $\mu_{\text{calm}}$ column and the collapsed component (always at $\mu = -27.63$) in the $\mu_{\text{crisis}}$ column; this matches the post-fallback labelling used downstream (top 20\% by log-vol = crisis), so the calm / crisis assignment in this table is bit-stable with the assignment used by the ARI computations regardless of degeneracy. Separation = $|\mu_{\text{crisis}} - \mu_{\text{calm}}| / \sigma_{\text{pooled}}$. Overlap = sum of tail probabilities beyond the midpoint boundary. Cells marked \textbf{degen.}: one GMM component collapses to near-zero variance ($\sigma \approx 0.001$, weight $< 0.005$), so the raw GMM assigns $\approx$100\% of bars to the higher-mean component; the percentile fallback (top 20\% = crisis) is triggered, and the ARI values in Tables~1 and \ref{tab:ari_matrix_combined} use the post-fallback labels. \emph{USD/JPY sample-scale note:} the table reveals a large disparity in $n$ across frequencies (36,393 5m bars vs.\ 156 1d bars in 2026). This reflects the tri-session structure (Tokyo $\approx$21:00--06:00, London $\approx$03:00--12:00, New York $\approx$08:00--17:00 in Eastern time), each contributing roughly one-third of the 5m sample; the 1d GMM therefore fits on 156 calendar-day observations, while ARI between forward-filled 1d labels and 5m labels is computed on all 36,393 5m bars. The 1h fit ($n=3{,}067$, separation 1.13) anchoring the 15m$\to$1h structural-break comparison is well-populated.
\par}
\end{table}

\paragraph{Expanding-window estimation}
Table~\ref{tab:expanding} compares full-sample and expanding-window ARI. Expanding-window ARI is lower than the full-sample ARI for all five 2026 assets (SPY $-0.066$, QQQ $-0.022$, CL $-0.284$, USD/JPY $-0.020$, GLD $-0.062$), consistent with early-window estimation noise being absorbed by the full-sample fit; the five-asset expanding-window mean of 0.108 lies essentially at the calibrated $K{=}2$ MS-Gaussian baseline (0.113, IQR 0.099--0.125 from canonical-pipeline simulation; \S\ref{app:rss_sim}). The empirical-vs-calibrated alignment without look-ahead is preserved for SPY, QQQ, USD/JPY and GLD; CL's expanding-window estimate (0.110) is materially below its full-sample reading (0.394) because the expanding warm-up samples a span before the 28 February 2026 supply-shock cluster that drives CL's headline 1h--1d agreement.

% source: outputs/pipeline_summary.csv  (overall_mean_ari + expanding_mean_ari, indexed by symbol)
\begin{table}[H]
\centering
\small
\setlength{\tabcolsep}{5pt}
\caption{Expanding-window GMM: mean off-diagonal ARI vs.\ full-sample baseline.}
\label{tab:expanding}
\begin{tabular}{lcc}
\toprule
Asset & Full-sample ARI & Expanding-window ARI \\
\midrule
SPY & 0.183 & 0.117 \\
QQQ & 0.157 & 0.135 \\
CL & 0.394 & 0.110 \\
USD/JPY & 0.105 & 0.085 \\
GLD & 0.155 & 0.093 \\
\bottomrule
\end{tabular}
\par\medskip
{\footnotesize\justifying
\textit{Notes:} At each bar $t$, GMM fitted on $[1, t]$; label is the out-of-sample prediction. Refitted every $n/50$ bars for computational efficiency.
\par}
\end{table}

\paragraph{Block-permutation test}
Table~\ref{tab:block_perm} reports block-permutation results.

% source: outputs/hypothesis_tests.csv  (filter test_type='block_permutation', block_size=50)
\begin{table}[H]
\centering
\small
\setlength{\tabcolsep}{5pt}
\caption{Block-permutation test for mean off-diagonal ARI (block size = 50 bars).}
\label{tab:block_perm}
\begin{tabular}{lccc}
\toprule
Asset & Off-diag.\ ARI & $p$-value & Null 95\% CI \\
\midrule
SPY & 0.183 & $<$0.002 & [$-$0.001, 0.004] \\
QQQ & 0.157 & $<$0.002 & [$-$0.003, 0.006] \\
CL & 0.394 & $<$0.002 & [$-$0.008, 0.010] \\
USD/JPY & 0.105 & $<$0.002 & [$-$0.007, 0.008] \\
GLD & 0.155 & $<$0.002 & [$-$0.012, 0.015] \\
\bottomrule
\end{tabular}
\par\medskip
{\footnotesize\justifying
\textit{Notes:} Block size = 50 bars ($\approx$ 4 hours at 5m). Within-block autocorrelation preserved; cross-resolution temporal alignment destroyed. 500 permutations; $p < 0.002$ ($= 1/501$ with add-one smoothing, as in Table~\ref{tab:perm_overall}). Null CI is wider than in Table~\ref{tab:perm_overall} because regime persistence is retained.
\par}
\end{table}

\paragraph{Calendar-window-length sensitivity}
Table~\ref{tab:window_sweep} reports how the intraday-only mean ARI varies with the common rolling-volatility window length.

% source: outputs/{SPY,CL,USDJPY,GLD}_calendar_window_cross_freq_ari.csv  (intraday-only mean per window length)
\begin{table}[H]
\centering
\small
\setlength{\tabcolsep}{5pt}
\caption{Calendar-window sweep: intraday-only mean off-diagonal ARI by window length.}
\label{tab:window_sweep}
\begin{tabular}{lccccc}
\toprule
Asset & 2h & 4h & 6h & 12h & 24h \\
\midrule
SPY & 0.253 & 0.326 & 0.369 & 0.555 & 0.039 \\
QQQ & --- & --- & 0.367 & --- & --- \\
CL & 0.216 & 0.615 & 0.670 & 0.407 & 0.444 \\
USD/JPY & 0.128 & 0.049 & 0.144 & 0.582 & 0.675 \\
GLD & 0.216 & 0.257 & 0.493 & 0.661 & 0.843 \\
\bottomrule
\end{tabular}
\par\medskip
{\footnotesize\justifying
\textit{Notes:} Each column applies the same fixed rolling-volatility window to all three intraday frequencies (5m, 15m, 1h); daily is excluded. QQQ 6h is sourced from \texttt{outputs/pipeline\_summary.csv}; the 2h/4h/12h/24h values for QQQ are not persisted in \texttt{outputs/} and are shown as ---. \emph{Window-length effects are asset-specific and session-dependent, not collapsing to a single sweet spot.} The five-asset cross-asset band at 6h is [0.14, 0.67] (width 0.53); at 12h the four verified assets span [0.41, 0.66] (width 0.25), so 12h is more consistent across assets than 6h. The 6h primary specification is selected for being non-pathological across all five assets (no asset crashes to ${<}\,0.10$ or rises above 0.7) rather than for cross-asset consistency. The asset-specific window dependence is itself evidence of session-driven dissonance: \SPY{} peaks at 12h (within its $\sim 13$h extended trading session) and collapses to 0.04 at 24h because the 24-hour window spans the extended-session boundary; \USDJPY{} rises monotonically through 24h matching its 24-hour FX market structure; \CL{} peaks at 4h--6h matching the NY-session arc of supply-shock realisation; \GLD{} rises monotonically through 24h matching its global safe-haven-flow window. These patterns reinforce the microstructure framing of \S{}3.2 (session-structure heterogeneity drives cross-asset inversion).
\par}
\end{table}

\paragraph{Block-size sensitivity}
Table~\ref{tab:block_sweep} shows that the block-permutation conclusion is insensitive to block-size choice.

% source: outputs/block_sweep_assets.csv  (SPY/QQQ/CL/USDJPY rows) + outputs/gld_block_sweep.csv (GLD row)
\begin{table}[H]
\centering
\small
\setlength{\tabcolsep}{5pt}
\caption{Block-permutation $p$-value and null 95\% upper bound by block size.}
\label{tab:block_sweep}
\begin{tabular}{lcccc}
\toprule
 & \multicolumn{4}{c}{Block size (5m bars)} \\
\cmidrule(lr){2-5}
Asset & 25 & 50 & 100 & 250 \\
\midrule
SPY $p$ & $<$0.002 & $<$0.002 & $<$0.002 & $<$0.002 \\
\quad null 95\% UB & 0.003 & 0.004 & 0.005 & 0.008 \\
QQQ $p$ & $<$0.002 & $<$0.002 & $<$0.002 & $<$0.002 \\
\quad null 95\% UB & 0.004 & 0.006 & 0.007 & 0.011 \\
CL $p$ & $<$0.002 & $<$0.002 & $<$0.002 & $<$0.002 \\
\quad null 95\% UB & 0.008 & 0.010 & 0.014 & 0.021 \\
USD/JPY $p$ & $<$0.002 & $<$0.002 & $<$0.002 & $<$0.002 \\
\quad null 95\% UB & 0.007 & 0.008 & 0.011 & 0.013 \\
GLD $p$ & $<$0.002 & $<$0.002 & $<$0.002 & $<$0.002 \\
\quad null 95\% UB & 0.013 & 0.015 & 0.020 & 0.028 \\
\bottomrule
\end{tabular}
\par\medskip
{\footnotesize\justifying
\textit{Notes:} 500 permutations per setting; observed ARI is the full-sample mean off-diagonal (SPY 0.184, QQQ 0.157, CL 0.392, USD/JPY 0.105, GLD 0.154). The null 95\% upper bound grows monotonically with block size (as expected: larger blocks preserve more autocorrelation), but remains $5.5$--$63\times$ below the observed ARI at every block size; the smallest margin is GLD at the 250-bar block ($5.5\times$). All five 2026 assets reject the regime-permuted null at $p < 0.002$ at every block size, including the autocorrelation-conservative 250-bar block ($\approx$1 trading day).
\par}
\end{table}

\paragraph{CL roll-week analysis}
Table~\ref{tab:cl_roll} splits the CL sample at contract-roll weeks; the negligible difference rules out roll contamination.

% source: outputs/CL_roll_week_ari.csv
\begin{table}[H]
\centering
\small
\setlength{\tabcolsep}{5pt}
\caption{CL (WTI continuous front-month): ARI by roll-week vs.\ non-roll-week periods.}
\label{tab:cl_roll}
\begin{tabular}{lcc}
\toprule
Period & Mean off-diag.\ ARI & 5m bars \\
\midrule
Roll weeks (days 18--22) & 0.388 & 5,562 \\
Non-roll weeks & 0.392 & 29,334 \\
\bottomrule
\end{tabular}
\par\medskip
{\footnotesize\justifying
\textit{Notes:} Roll weeks approximated as calendar days 18--22 of each month (the core of the WTI front-month roll period), matching the code heuristic in \texttt{cl\_roll\_week\_analysis}. Within each roll week, the small subset of bars carrying contract-rollover return jumps (identified by an asset-specific MAD-based filter, \S{}2.2) is replaced with median-Winsorised values before volatility estimation; the remaining bars in the roll-week subset are retained for ARI computation. The 0.388 figure is therefore computed on roll-week bars under cleaned-jump returns rather than on raw rollover returns, and the 5{,}562-bar count reflects all retained 5m bars in roll weeks (not the small subset of jump-affected bars).
\par}
\end{table}

\paragraph{Out-of-sample validation: 2022 Russia--Ukraine}
Table~\ref{tab:oos_2022} reports the 2022 out-of-sample results.

% source: outputs_2022/pipeline_summary.csv                       (Baseline, Expand. window)
%         outputs_2022/{ASSET}_tod_adjusted_cross_freq_ari.csv    (TOD-adj.)
%         outputs_2022/{ASSET}_calendar_window_cross_freq_ari.csv (Cal.-aligned, intraday)
%         outputs_2022/{ASSET}_cross_freq_ari.csv                 (pairwise: 5m-15m, 5m-1h, 1h-1d)
%         p_block recomputed inline by src/experiments/exp_03_hypothesis_tests.py for 2022 outputs
\begin{table}[H]
\centering
\small
\setlength{\tabcolsep}{4pt}
\caption{Out-of-sample validation: 2022 Russia--Ukraine episode (December 2021 -- May 2022). \CL{} omitted because IB \texttt{CONTFUT} historical service does not retain pre-2024 5m data; \GLD{} provides commodity-class coverage.}
\label{tab:oos_2022}
\resizebox{\textwidth}{!}{%
\begin{tabular}{lcccccccc}
\toprule
 & \multicolumn{4}{c}{Mean off-diag.\ ARI} & & \multicolumn{3}{c}{Pairwise ARI} \\
\cmidrule(lr){2-5}\cmidrule(lr){7-9}
Asset & Baseline & Expand.\ window & TOD-adj. & Cal.-aligned (intra.) & $p_{\text{block}}$ & 5m--15m & 5m--1h & 1h--1d \\
\midrule
SPY & 0.146 & 0.133 & 0.329 & 0.360 & 0.002 & 0.612 & 0.159 & 0.007 \\
USD/JPY & 0.153 & 0.097 & 0.362 & 0.450 & 0.002 & 0.549 & 0.190 & $-$0.006 \\
GLD & 0.133 & 0.060 & 0.302 & 0.187 & 0.002 & 0.656 & 0.139 & $-$0.032 \\
\bottomrule
\end{tabular}%
}
\par\medskip
{\footnotesize\justifying
\textit{Notes:} Same methodology as the 2026 primary sample. Expand.\ window = expanding-window GMM without look-ahead. $p_{\text{block}}$ = block-permutation $p$-value (block size = 50). Stress window = 22--28 February 2022; calm window = 10--14 January 2022 (both 5 trading days).
\par}
\end{table}

\paragraph{Majority-vote upward aggregation}
Table~\ref{tab:majority_vote} reports majority-vote ARI per asset and episode.

% source: outputs/majority_vote_summary.csv       (2026)
%         outputs_2022/majority_vote_summary.csv  (2022 OOS)
\begin{table}[H]
\centering
\small
\setlength{\tabcolsep}{6pt}
\caption{Majority-vote upward aggregation: per-asset \emph{six-pair} mean off-diagonal ARI on the coarse grid. Each coarse label = mode of the constituent 5m labels (ties resolve to crisis); ARI is then computed pair-wise across all six $4 \times 4$ off-diagonal cells (5m--15m, 5m--1h, 5m--1d, 15m--1h, 15m--1d, 1h--1d) on the coarse grid and the per-asset mean is reported. Adjacent-pair entries (5m--15m, 15m--1h, 1h--1d) are also reported in Table~1 for direct comparison with the forward-fill column.}
\label{tab:majority_vote}
\begin{tabular}{lccc}
\toprule
Asset & 2026 ARI & 2022 ARI & Combined range \\
\midrule
SPY & 0.231 & 0.127 & 0.127--0.231 \\
QQQ & 0.209 & --- & --- \\
CL & 0.421 & --- & --- \\
USD/JPY & 0.096 & 0.155 & 0.096--0.155 \\
GLD & 0.148 & 0.134 & 0.134--0.148 \\
\bottomrule
\end{tabular}
\par\medskip
{\footnotesize\justifying
\textit{Notes:} Both 5m and the coarser frequency are independently fitted by GMM; we then aggregate the 5m labels to the coarser grid by per-bin majority vote (ties at exactly 50\% crisis-share resolved to crisis: $\text{share} \ge 0.5 \mapsto 1$) and compute ARI against the coarser GMM labels at the coarse grid. \CL{} 2022 not available (IB CONTFUT historical service does not retain pre-2024 5m data). The combined range across both episodes is 0.09--0.42, comparable to the forward-fill baseline (Table~1) and well below the calendar-aligned 6h intraday ceiling at the asset level (Table~\ref{tab:decomposition}).
\par}
\end{table}

\paragraph{Bootstrap test over 5-day windows}
Table~\ref{tab:bootstrap_5d} reports 1{,}000-window bootstrap statistics for the mean off-diagonal ARI, with two-sided $p$-values for the calm and stress windows relative to the bootstrap distribution.

% source: outputs/bootstrap_five_day_windows.csv       (2026)
%         outputs_2022/bootstrap_five_day_windows.csv  (2022 OOS)
\begin{table}[H]
\centering
\small
\setlength{\tabcolsep}{4pt}
\caption{Bootstrap test: mean off-diag.\ ARI over 1{,}000 random 5-day windows.}
\label{tab:bootstrap_5d}
\resizebox{\textwidth}{!}{%
\begin{tabular}{llccccccc}
\toprule
Episode & Asset & Calm ARI & Stress ARI & Boot.\ median & Boot.\ 95\% interval & $p_{\text{calm}}$ & $p_{\text{stress}}$ \\
\midrule
2026 & SPY & 0.097 & 0.202 & 0.115 & [0.054, 0.270] & 0.33 & 0.08 \\
2026 & QQQ & 0.082 & 0.142 & 0.106 & [0.066, 0.267] & 0.36 & 0.28 \\
2026 & USD/JPY & 0.079 & 0.052 & 0.081 & [$-$0.006, 0.455] & 0.73 & 0.93 \\
2026 & CL & 0.151 & 0.079 & 0.205 & [0.059, 1.000] & 0.99 & 0.98 \\
2026 & GLD & 0.125 & 0.141 & 0.133 & [0.044, 0.301] & 0.62 & 0.22 \\
2022 & SPY & 0.104 & 0.160 & 0.114 & [0.076, 0.263] & 0.76 & 0.26 \\
2022 & USD/JPY & 0.094 & 0.103 & 0.110 & [0.063, 0.286] & 0.75 & 0.56 \\
2022 & GLD & 0.050 & 0.128 & 0.114 & [0.014, 0.294] & 0.24 & 0.68 \\
\bottomrule
\end{tabular}%
}
\par\medskip
{\footnotesize\justifying
\textit{Notes:} Bootstrap windows are 5 contiguous trading days drawn uniformly from the asset's calendar. GMM regimes are fitted once on the full sample to isolate sampling variability of ARI from refit instability. The mean off-diagonal ARI within each 5-day window is computed over the three intraday-intraday pairs only (5m--15m, 5m--1h, 15m--1h); pairs involving 1d return NaN at this window length because the 1d series has fewer than the 10-bar minimum required for ARI estimation. The reported statistic is therefore the intraday-only 3-pair mean, distinct from the 6-pair mean in Table~1. $p_{\text{calm}}, p_{\text{stress}}$ are two-sided fractions of bootstrap statistics whose deviation from the bootstrap median exceeds the calm or stress ARI's deviation; we read these as descriptive tail-rank statistics of the calm and stress windows within the bootstrap distribution rather than as formal hypothesis tests. The 16 calm/stress vs.\ bootstrap comparisons across the eight asset--episode combinations satisfy $p \in [0.08, 0.99]$ with median $p = 0.59$; five values fall below the conventional $p > 0.3$ threshold (\SPY{} 2026 stress at 0.08 is the lowest; QQQ 2026 stress 0.28, \GLD{} 2026 stress 0.22, \SPY{} 2022 stress 0.26, \GLD{} 2022 calm 0.24 are the others). The calm and stress windows are therefore not systematically unusual realisations relative to a random 5-day draw from the same calendar, with the exception of \SPY{} 2026 stress, which occupies the upper tail of the bootstrap distribution.
\par}
\end{table}

\paragraph{Resolution-conditional VaR uplift on disagreement days}
Table~\ref{tab:var_uplift} reports the population statistics of the regime-conditional 99\% empirical-VaR ratio underlying the conservative-resolution VaR statistic of \S4. The statistic is reported \emph{ex post} on full-sample quantiles and is not a realisable posting rule; a causal expanding-window analogue is left for follow-on work. For each asset and episode we fit GMMs at every frequency on the full sample, compute the 99\% empirical quantile of 5-minute log-returns within each regime under the 1h and 1d classifiers separately, and aggregate to per-day labels using the \emph{native}-resolution label series (1d: one bar per trading day; 1h: within-day majority over the 1h labels with ties resolved to crisis, $\text{share} \ge 0.5 \mapsto 1$). Native-resolution aggregation is the operative protocol because the alternative (forward-filling 1d labels onto the 5m grid and then taking a calendar-day mode-vote) assigns the previous trading day's 1d label to bars before 16:00 NY, producing a 1-day phase-shifted day-level signal that distorts the disagreement-day count. The disagreement set is the subset of days where the 1h-classifier and 1d-classifier labels differ; on those days the ratio statistic is $\max(|\text{VaR}_{1h}|, |\text{VaR}_{1d}|)/\min(|\text{VaR}_{1h}|, |\text{VaR}_{1d}|)$. The full-sample uplift column reports the average percentage increase in posted VaR if the conservative resolution is always adopted (relative to a 1d-only baseline). \CL{} is reported only for 2026 because IB \texttt{CONTFUT}'s pre-2024 5m history is forward-fill placeholder data (\S{}2.1).

% source: outputs/var_uplift_by_resolution.csv       (2026)
%         outputs_2022/var_uplift_by_resolution.csv  (2022 OOS)
\begin{table}[H]
\centering
\small
\setlength{\tabcolsep}{4pt}
\caption{Resolution-conditional 99\% empirical-VaR uplift on 1h-vs-1d disagreement days.}
\label{tab:var_uplift}
\resizebox{\textwidth}{!}{%
\begin{tabular}{llcccccc}
\toprule
Episode & Asset & Days & Disagree (\%) & Ratio median & Ratio q75 & Avg.\ uplift (full) & Trustworthy \\
\midrule
2026 & SPY & 122 & 26.2 & 1.33 & 1.33 & 7.5\% & yes \\
2026 & QQQ & 122 & 27.0 & 1.42 & 1.42 & 4.2\% & yes \\
2026 & USD/JPY & 154 & 52.6 & 1.16 & 1.16 & 7.8\% & yes \\
2026 & CL & 153 & 17.6 & 2.24 & 2.24 & 21.0\% & yes \\
2026 & GLD & 122 & 63.1 & 1.51 & 1.51 & 13.9\% & yes \\
2022 & SPY & 123 & 43.9 & 1.23 & 1.23 & 10.0\% & yes \\
2022 & USD/JPY & 154 & 41.6 & 1.46 & 1.46 & 12.3\% & yes \\
2022 & GLD & 123 & 49.6 & 1.27 & 1.27 & 12.7\% & yes \\
\bottomrule
\end{tabular}%
}
\par\medskip
{\footnotesize\justifying
\textit{Notes:} 99\% empirical VaR computed on 5m log-returns pooled by regime under each classifier. Full-sample uplift formula: $\bar{u} = T^{-1}\sum_t (\max(|\text{VaR}^{1h}_t|,|\text{VaR}^{1d}_t|)/|\text{VaR}^{1d}_t|) - 1$ averaged across all valid days $t$.

\emph{Single-value framing.} The regime-conditional VaR is one of four configurations and is computed once per (asset, configuration) cell from full-sample empirical quantiles, so the disagreement-day ratio takes at most two distinct values across days; median = q75 reflects dominant-configuration concentration, not statistical robustness. Dominant-configuration shares among disagreement days (\texttt{outputs/disagree\_config\_breakdown.csv}): SPY 2026 87.5\% 1h-crisis/1d-calm (28/32); QQQ 2026 87.9\% 1h-calm/1d-crisis (29/33); USD/JPY 2026 97.5\% 1h-calm/1d-crisis (79/81); CL 2026 96.3\% 1h-crisis/1d-calm (26/27); GLD 2026 100.0\% 1h-calm/1d-crisis (77/77); SPY 2022 100.0\% 1h-crisis/1d-calm (54/54); USD/JPY 2022 90.6\% 1h-calm/1d-crisis (58/64); GLD 2022 90.2\% 1h-calm/1d-crisis (55/61). The economically meaningful statistic is the magnitude of the configuration-pair gap (``if 1h and 1d disagree, the resulting VaRs differ by this factor''); day-level uncertainty enters through the disagreement frequency column (18--63\% across asset--episodes), not through ratio dispersion.

\emph{Decomposition.} The full-sample uplift has two channels: a \emph{disagreement-day} channel ($p_d$, mean uplift $u_d$) where the $\max$-rule may or may not bind depending on which classifier flags the day as crisis, and an \emph{agreement-day} channel ($p_a = 1 - p_d$, mean uplift $u_a$) where both classifiers label the day identically but the regime-conditional VaR pools differ in composition (1h and 1d classifiers can disagree on which 5m bars belong in the ``crisis pool''). The check: $\bar{u} = p_d u_d + p_a u_a$. CL 2026's headline 21.0\% reflects both channels: the dominant 1h-crisis/1d-calm configuration on disagreement days lifts $|\text{VaR}^{1h}|$ above the 1d crisis baseline, while pool-composition asymmetry between 1h and 1d crisis pools (1h crisis 13{,}105 5m bars vs.\ 1d crisis 7{,}437; i.e.\ the 1d classifier flags fewer bars as crisis even on agreement-crisis days) loads on label-agreement-crisis days. Per-regime pool counts are in \texttt{outputs/var\_uplift\_by\_resolution.csv}; the table above carries only the dominant-configuration ratio for headline parsimony.

All eight (asset, episode) rows pass the trustworthy gate (each regime has ${\ge}\,50$ 5m bars; no degenerate flag). CL is reported only for 2026 because IB's CONTFUT historical service does not retain pre-2024 front-month 5m data; QQQ 2022 5m bars were not pulled in this revision round; GLD provides commodity-class coverage in both episodes.
\par}
\end{table}

\paragraph{Calm-day-only subsample}
Table~\ref{tab:calm_day_subsample} reports the calm-day-only subsample analysis described in \S{}3.4: days with daily realised volatility below the \emph{per-asset in-sample median} and outside the peak-stress window are retained, and the cross-frequency ARI is recomputed on this subset using the full-sample GMM boundaries. Per-asset (rather than cross-asset) cutoffs preserve each asset's relative quiet-day structure but mean that, e.g., \CL{}'s calm-day subsample contains days that would lie above SPY's stress threshold in absolute realised-vol terms.

% source: outputs/calm_day_subsample_ari.csv       (2026)
%         outputs_2022/calm_day_subsample_ari.csv  (2022 OOS)
\begin{table}[H]
\centering
\small
\setlength{\tabcolsep}{4pt}
\caption{Calm-day-only subsample: full-sample vs.\ calm-only mean off-diag.\ ARI and pairwise structure.}
\label{tab:calm_day_subsample}
\resizebox{\textwidth}{!}{%
\begin{tabular}{llccccccc}
\toprule
Episode & Asset & Calm days $n$ & Total days & Full ARI & Calm ARI & 5m--15m & 15m--1h & 1h--1d \\
\midrule
2026 & SPY & 54 & 124 & 0.183 & 0.110 & 0.461 & 0.050 & 0.072 \\
2026 & QQQ & 51 & 124 & 0.157 & 0.100 & 0.496 & 0.067 & $-$0.034 \\
2026 & USD/JPY & 54 & 156 & 0.105 & 0.051 & 0.256 & 0.049 & 0.010 \\
2026 & CL & 76 & 155 & 0.394 & 0.080 & 0.355 & 0.065 & 0.000 \\
2026 & GLD & 52 & 124 & 0.155 & 0.107 & 0.577 & 0.028 & $-$0.008 \\
2022 & SPY & 62 & 125 & 0.146 & 0.110 & 0.541 & 0.051 & 0.000 \\
2022 & USD/JPY & 74 & 156 & 0.153 & 0.123 & 0.531 & 0.130 & $-$0.026 \\
2022 & GLD & 60 & 125 & 0.133 & 0.110 & 0.556 & 0.070 & $-$0.003 \\
\bottomrule
\end{tabular}%
}
\par\medskip
{\footnotesize\justifying
\textit{Notes:} Calm-day subsample = days with daily realised volatility (root-sum-of-squares of 5m log-returns) strictly below the in-sample median, additionally excluding all dates in the peak-stress window (2026-02-28 to 2026-04-15; 2022-02-22 to 2022-02-28). GMM regimes fitted on the full sample; ARI recomputed on the calm-day 5m timestamps. The hierarchical pattern is preserved on calm-only days: high 5m--15m agreement (0.26--0.56 in 2026, 0.53--0.55 in 2022) collapses across the 15m$\to$1h boundary to $|\text{ARI}|\le 0.13$, and the calm-subsample 1h--1d pair is uniformly low across all eight asset--episode rows ($|\text{ARI}|\le 0.08$, including SPY 2026 calm at 0.072). The full-sample 1h--1d pair is similarly near zero for CL, USD/JPY, and GLD but reaches 0.171 for SPY 2026 (Table~1); the calm-subsample comparison strips out the disagreement-day component of SPY's 1h--1d signal, leaving only the near-zero residual that the other assets show across the full sample. Resolution dissonance is thus not a stress-period artefact.
\par}
\end{table}

\paragraph{Paired test: stress-window vs calm-window ARI}
Table~\ref{tab:stress_vs_calm} formalises the calm-vs-stress comparison referenced in \S{}3.4. For each of the eight asset--episode combinations (five 2026 assets plus three 2022 OOS assets; \CL{} 2022 absent because IB CONTFUT pre-2024 5m bars are placeholders; QQQ 2022 not pulled in this round) we take the stress-window and calm-window ARI already reported in Table~\ref{tab:bootstrap_5d}, and test the paired difference.

% source: outputs/stress_vs_calm_test.csv          (per-asset paired diffs)
%         outputs/stress_vs_calm_test_summary.csv  (paired t / Wilcoxon / sign p-values)
\begin{table}[H]
\centering
\small
\setlength{\tabcolsep}{5pt}
\caption{Paired test of stress-window vs calm-window mean off-diag.\ ARI. All hypothesis tests fail to reject equality.}
\label{tab:stress_vs_calm}
\begin{tabular}{llccc}
\toprule
Episode & Asset & Calm ARI & Stress ARI & Diff (stress $-$ calm) \\
\midrule
2026 US--Iran & SPY & 0.097 & 0.202 & $+$0.105 \\
2026 US--Iran & QQQ & 0.082 & 0.142 & $+$0.060 \\
2026 US--Iran & USD/JPY & 0.079 & 0.052 & $-$0.027 \\
2026 US--Iran & CL & 0.151 & 0.079 & $-$0.072 \\
2026 US--Iran & GLD & 0.125 & 0.141 & $+$0.015 \\
2022 Russia--Ukraine & SPY & 0.104 & 0.160 & $+$0.055 \\
2022 Russia--Ukraine & USD/JPY & 0.094 & 0.103 & $+$0.009 \\
2022 Russia--Ukraine & GLD & 0.050 & 0.128 & $+$0.077 \\
\midrule
\multicolumn{3}{l}{Mean difference} & & $+$0.028 \\
\multicolumn{3}{l}{Std.\ deviation of differences} & & 0.058 \\
\multicolumn{3}{l}{Paired $t$-test} & $t = +1.36$ & $p = 0.216$ \\
\multicolumn{3}{l}{Wilcoxon signed-rank} & $W = 9.0$ & $p = 0.250$ \\
\multicolumn{3}{l}{Exact sign test} &  & $p = 0.289$ \\
\multicolumn{3}{l}{Stress $>$ calm} & 6 / 8 pairs & \\
\bottomrule
\end{tabular}
\par\medskip
{\footnotesize\justifying
\textit{Notes:} 2026 calm window: 2026-01-01 to 2026-01-24 (pre-escalation baseline per \texttt{EPISODES['2026\_iran']}); stress window: 2026-02-28 to 2026-04-15 (US--Israeli strike through early-April ceasefire). 2022 calm window: 2022-01-10 to 2022-01-14; stress window: 2022-02-22 to 2022-02-28. ARI values are the full-episode-window mean off-diagonal ARI (column \texttt{calm\_ari} / \texttt{stress\_ari} of \texttt{bootstrap\_five\_day\_windows.csv}), not 5-day-window ARIs. All three hypothesis tests fail to reject $H_0:$ stress ARI $=$ calm ARI at any conventional significance level. Combined with the bootstrap $p$-values in Table~\ref{tab:bootstrap_5d}, this is consistent with cross-frequency label disagreement \emph{not} being a stress-period artefact. \emph{Cross-episode directional pattern (post-hoc, not statistically significant):} the 6/8 stress$>$calm split is asymmetric across episodes: 3/3 in 2022 OOS (all positive: $+0.055$, $+0.009$, $+0.077$) versus 3/5 in 2026 primary (SPY $+0.105$, QQQ $+0.060$, GLD $+0.015$ positive; USD/JPY and CL negative). The two equity ETFs (SPY, QQQ) both show stress-elevation; FX and futures show mixed signs. Neither pattern is significant at conventional levels and we do not interpret this asymmetry as a finding, but report it for transparency.
\par}
\end{table}

\subsection{Kruskal--Wallis test for hierarchical ARI structure}\label{app:kw_test}

Pooling across the eight asset--episode combinations (five 2026 + three 2022 OOS), mean ARI is 0.589 for adjacent intraday pairs (5m--15m, $n=8$), 0.177 for non-adjacent intraday pairs (5m--1h, 15m--1h; $n=16$), and 0.042 for intraday--daily pairs (5m--1d, 15m--1d, 1h--1d; $n=24$). The pooled Kruskal--Wallis test rejects equality of medians ($H = 29.57$, $p \approx 3.8 \times 10^{-7}$); the per-episode tests are individually significant ($H_{2026}=16.07$, $p_{2026}=3.2\times 10^{-4}$; $H_{2022}=14.21$, $p_{2022}=8.2\times 10^{-4}$) and combine via Fisher's method to $\chi^2_{\text{Fisher}}=30.28$, $p_{\text{Fisher}}=4.3\times 10^{-6}$ (the recommended omnibus statistic since each episode contributes independent draws). Mann--Whitney $U$-tests confirm both gaps: adjacent vs.\ non-adjacent intraday $p \approx 1.4 \times 10^{-6}$, non-adjacent intraday vs.\ intraday--daily $p \approx 8.2 \times 10^{-5}$ (BH-FDR adjusted). The tiered hierarchy is therefore evident at both 5m--15m vs.\ wider intraday and intraday vs.\ intraday--daily transitions. \emph{Grouping disclosure (post-hoc, not pre-registered).} The three-group partition reflects the visual tiering observed in \S{}3.1: adjacent fine-pair (5m--15m) at the high sub-band, non-adjacent intraday (5m--1h, 15m--1h) spanning the 15m$\to$1h single-step drop, and intraday--daily (5m--1d, 15m--1d, 1h--1d) at the cross-class horizon. The grouping is data-driven from the pairwise inspection in \S{}3.1 rather than pre-registered; the test should therefore be read as a quantification of the visual pattern, not as confirmatory hypothesis testing. The pooled $H$ statistic uses the asymptotic chi-square approximation on $n_1 = 8$, $n_2 = 16$, $n_3 = 24$, while the per-episode tests use $n=5, 10, 15$ for 2026 and $n=3, 6, 9$ for 2022; the Fisher's-method combination $p_{\text{Fisher}}=4.3\times 10^{-6}$ is the preferred summary because it respects the within-asset cross-episode correlation that biases the pooled $p$ low.

\subsection{Simulation-based calibration: benchmarking empirical ARI against theory}\label{app:rss_sim}

\citet{ChanZhangCheung2009} predict that aggregation transforms the DGP, so \emph{some} cross-frequency label disagreement is a structural certainty. We construct a calibrated 2-state Markov-switching Gaussian baseline under apples-to-apples comparison: synthetic 5-minute returns are generated from the calibrated DGP, built into a synthetic OHLC frame, and run through the same canonical analysis pipeline (\texttt{fit\_regimes\_per\_frequency} + \texttt{cross\_freq\_ari\_matrix}) used on the empirical data. The baseline mean off-diagonal ARI from $n_{\text{rep}}=200$ replications is the directly comparable reference for the empirical ARI of 0.11--0.39 reported in Table~1.

\paragraph{Data-generating process}
DGP parameters are ML-fit on SPY 1h log returns over the empirical sample (Nov 2025--May 2026, $n=2{,}097$ 1h bars) using \texttt{statsmodels.tsa.regime\_switching.MarkovRegression} with $K=2$ and switching variance. The fitted persistences are $P_{12}=0.197$ and $P_{21}=0.256$ at 1h scale, corresponding to mean regime durations of $\approx 5$ and $\approx 4$ hours respectively. State means and standard deviations are sorted so the lower-variance state is calm and the higher-variance state is crisis. The 5m-scale parameters used to drive the simulator are obtained as the principal $1/12$-th matrix root of the $2\times 2$ 1h transition matrix via \texttt{scipy.linalg.fractional\_matrix\_power} (yielding $P_{12}=0.021$, $P_{21}=0.028$ at 5m, $\tau \equiv P_{12}+P_{21} = 0.049$). State-conditional moments scale as $\mu_{5m} = \mu_{1h}/12$, $\sigma_{5m} = \sigma_{1h}/\sqrt{12}$. Initial state at $t=0$ is drawn from the stationary distribution $\pi_0 = P_{21}/(P_{12}+P_{21})$. The simulator generates a 5m return path on a synthetic NY-tz RTH DatetimeIndex, builds an OHLC frame via $P_t = P_{t-1}(1+r_t)$ with high/low set to $\max/\min$ of close, resamples to 5m / 15m / 1h / 1d via the canonical \texttt{\_resample\_ohlc}, and runs \texttt{fit\_regimes\_per\_frequency} (the same fitter the empirical pipeline uses, with rolling-window features and the K-means-warmstart GMM) plus \texttt{cross\_freq\_ari\_matrix} on the resulting label dictionary. We report the mean off-diagonal ARI across the six frequency pairs (``all-4'') and across the three intraday pairs (5m, 15m, 1h; ``intra-3''). For sensitivity, we also report results at three persistence settings ($P_{12}=P_{21} \in \{0.005, 0.003, 0.001\}$ at 5m scale) chosen to span slower-transition DGPs; these settings are NOT calibrated to data.

\paragraph{Null hypothesis ($H_0$: no resolution effect)}
Under $H_0$ we independently permute each coarse-resolution label sequence before computing ARI, preserving marginal crisis shares while destroying temporal alignment. Across all settings and 200 replications, the null ARI concentrates near zero (all-4 mean $< 0.002$, 97.5th percentile $< 0.02$; intra-3 mean $\approx 0.000$, 97.5th percentile $< 0.004$), consistent with the empirical permutation baselines in Table~\ref{tab:perm_overall}.

\paragraph{Theoretical baseline ($H_1$: true DGP)}
Under the true DGP (no permutation), the baseline ARI from canonical-pipeline simulation is:

\begin{center}
\small
\begin{tabular}{lcccc}
\toprule
$P_{12}$ & Duration (5m) & All-4 ARI [IQR] & Intra-3 ARI [IQR] \\
\midrule
calibrated$^\dagger$ & $\approx$47 / 36 bars & 0.113\; [0.099, 0.125] & 0.195\; [0.177, 0.213] \\
0.005 & $\approx$200 bars & 0.207\; [0.169, 0.232] & 0.329\; [\,---\,] \\
0.003 & $\approx$330 bars & 0.283\; [0.230, 0.331] & 0.411\; [\,---\,] \\
0.001 & $\approx$1{,}000 bars & 0.476\; [0.395, 0.564] & 0.603\; [\,---\,] \\
\bottomrule
\end{tabular}
\end{center}
\noindent $^\dagger$\emph{Calibrated setting}: $P_{12}=0.021$, $P_{21}=0.028$ at 5m, obtained as the principal $1/12$-th matrix root of the 1h ML fit on SPY 1h log returns ($P_{12}^{1h}=0.197$, $P_{21}^{1h}=0.256$; \S\ref{app:rss_sim} above). The remaining rows are sensitivity points at slower transition rates and are not calibrated.

\smallskip
\noindent At the calibrated DGP, the all-4 baseline ARI is 0.113 (IQR 0.099--0.125); the empirical 5-asset 2026 mean of 0.199 lies above this baseline (a +0.09 excess in mean, driven primarily by CL's $+0.27$ deviation; empirical readings for SPY/QQQ/USD/JPY/GLD cluster around the calibrated IQR while CL falls in the upper tail). The intra-3 baseline of 0.195 is exceeded by the empirical intraday-only 3-pair mean of 0.32 (per-asset 0.24--0.45 in 2026; 0.30--0.31 in 2022 OOS), giving an intraday excess of 0.05--0.26.

\paragraph{Robustness to richer DGP: GARCH(1,1)-Markov-switching baseline}
\label{app:garchms_sim}
To address the concern that the i.i.d.\ Gaussian innovation assumption may inflate the baseline by suppressing volatility clustering, we re-run the calibration replacing the Gaussian DGP with a regime-switching GARCH(1,1) process. GARCH parameters are ML-fit on SPY 5m log returns ($\alpha=0.318$, $\beta=0.578$, $\alpha+\beta=0.896$, stationary). Within regime $k$, $\sigma_t^2 = \omega_k + \alpha\,\epsilon_{t-1}^2 + \beta\,\sigma_{t-1}^2$, with regime-specific $\omega_k = \sigma_k^2(1-\alpha-\beta)$ calibrated so within-regime stationary variance matches the ML-fit MS-Gaussian variance, and on every regime change the GARCH state ($\sigma_{t-1}^2$, $\epsilon_{t-1}$) is reset to the new regime's stationary point so realised within-regime variance equals the calibrated $\sigma_k^2$ (matched-moments rather than upper-bound persistence). The same canonical-pipeline simulator drives the synthetic OHLC through \texttt{fit\_regimes\_per\_frequency} as in the Gaussian baseline. Results ($n_{\text{rep}}=200$):

\begin{center}
\small
\begin{tabular}{lccc}
\toprule
$P_{12}$ & DGP & All-4 ARI [IQR] & Intra-3 ARI \\
\midrule
\multirow{2}{*}{calibrated$^\dagger$}
 & Gaussian-MS  & 0.113\; [0.099, 0.125] & 0.195 \\
 & GARCH(1,1)-MS & 0.116\; [0.097, 0.131] & 0.202 \\[2pt]
\multirow{2}{*}{0.005}
 & Gaussian-MS  & 0.207\; [0.169, 0.232] & 0.329 \\
 & GARCH(1,1)-MS & 0.199\; [0.161, 0.236] & 0.305 \\[2pt]
\multirow{2}{*}{0.003}
 & Gaussian-MS  & 0.283\; [0.230, 0.331] & 0.411 \\
 & GARCH(1,1)-MS & 0.251\; [0.199, 0.298] & 0.365 \\[2pt]
\multirow{2}{*}{0.001}
 & Gaussian-MS  & 0.476\; [0.395, 0.564] & 0.603 \\
 & GARCH(1,1)-MS & 0.403\; [0.322, 0.495] & 0.523 \\
\bottomrule
\end{tabular}
\end{center}
\noindent $^\dagger$Calibrated row uses the ML-fit MS-Gaussian persistence at 5m ($P_{12}=0.021$, $P_{21}=0.028$, obtained as the principal $1/12$-th matrix root of the 1h fit).

\smallskip
\noindent At the calibrated DGP, the GARCH-MS baseline (0.116) is statistically indistinguishable from the Gaussian-MS baseline (0.113) at the 4-frequency level, and the intraday-only variants are within 0.01. Volatility clustering reduces the baseline ARI at slower persistence settings (where regime durations dominate), but the calibrated reference is essentially unchanged. The robustness check therefore rules out volatility clustering as a meaningful driver of the calibrated reference under apples-to-apples comparison. Extending the DGP to jump-diffusion or stochastic-volatility innovations is left for future work.

\emph{Microstructure-noise asymmetry between simulation and empirical pipelines.} The simulation generates noise-free 1-minute returns and aggregates to 5m, 15m, 1h, 1d; the empirical pipeline observes 5m bars that already contain bid--ask bounce and stale-quote contamination (\S\ref{app:data_choices}). The simulation 5m-fitting layer is therefore cleaner than the empirical 5m-fitting layer, and the calibrated baseline represents an upper bound on the noise-free ARI achievable under perfect fitting; introducing 5m i.i.d.\ noise (calibrated to Aït-Sahalia--Mykland two-scale estimates) at the simulation's 5m fitting layer would lower the baseline modestly, partly closing the empirical-baseline gap. The gap attributable to noise is bounded above by the AS-MK noise-variance-decay rate at 5m for the assets in our universe and is small relative to the gap to empirical (0.20 noise-corrected vs.\ 0.11--0.39 observed); a noise-corrected baseline simulation is left for future work. Separately, the simulation uses per-bar realised volatility (root-sum-of-squares of within-bar 1-minute returns) whereas the empirical pipeline uses a rolling-window estimator (\S{}2.2); the baseline ARI values are therefore indicative rather than exact calibration targets.

\subsection{$K$-matched calibrated baseline ($K{=}3$)}\label{app:k3_baseline}
The $K{=}2$ calibrated baseline above is the right reference for the $K{=}2$ empirical reading reported throughout the main text. For the $K{=}3$ robustness check (\S{}3.4) the fair reference is the $K{=}3$ canonical-pipeline baseline at the ML-fit persistence. The $K{=}3$ GMM is binarised by argmax-of-means (highest-mean state $=$ crisis) within the canonical pipeline, matching how $K{=}3$ is consumed in the empirical sweep. Slower-persistence sensitivity rows are reported for $K{=}2$ only because the $K{=}3$ DGP parameterisation at off-calibrated persistence grids does not have a clean three-state interpretation (see \texttt{exp\_13\_k3\_baseline.py}). Results ($n_{\text{rep}}=200$):

\begin{center}
\small
\begin{tabular}{lcccc}
\toprule
$P_{12}$ & $K$ & All-4 mean ARI [IQR] & Intraday mean ARI \\
\midrule
\multirow{2}{*}{calibrated$^\dagger$}
 & 2 & 0.113\; [0.099, 0.125] & 0.195 \\
 & 3 & 0.113\; [0.097, 0.125] & 0.196 \\[2pt]
\multirow{3}{*}{sensitivity ($K{=}2$ only)}
 & \multicolumn{3}{l}{$P_{12}=0.005$:\; 0.207\; [0.169, 0.232]\; intra 0.329} \\
 & \multicolumn{3}{l}{$P_{12}=0.003$:\; 0.283\; [0.230, 0.331]\; intra 0.411} \\
 & \multicolumn{3}{l}{$P_{12}=0.001$:\; 0.476\; [0.395, 0.564]\; intra 0.603} \\
\bottomrule
\end{tabular}
\end{center}
\noindent $^\dagger$Calibrated row uses the ML-fit MS-Gaussian persistence at 5m ($P_{12}=0.021$, $P_{21}=0.028$, obtained as the principal $1/12$-th matrix root of the 1h fit). $K{=}2$ sensitivity rows source: \texttt{outputs/simulation\_rss.csv}; $K{=}3$ calibrated row source: \texttt{outputs/simulation\_rss\_k3.csv}.

\smallskip
\noindent The $K{=}3$ calibrated baseline (0.113, IQR 0.097--0.125) is essentially identical to the $K{=}2$ calibrated baseline (0.113) at the canonical reference. The empirical $K{=}3$ five-asset 2026 mean (\texttt{outputs/robustness\_summary.csv} for $K{=}3$ rows) lies above the $K{=}3$ calibrated baseline, confirming the same intraday-coherence excess as the $K{=}2$ reading. \emph{Per-asset:} SPY 2026 (0.160), QQQ 2026 (0.140), USD/JPY 2026 (0.082), GLD 2026 (0.229), GLD 2022 (0.078), SPY 2022 (0.249), USD/JPY 2022 (0.241), and CL 2026 (0.390) sit at or above the $K{=}3$ calibrated reference (USD/JPY 2026 and GLD 2022 below), with CL 2026's 0.390 the largest excess (third state absorbs supply-shock days). The $K{=}3$ check therefore confirms that the empirical-vs-calibrated alignment is qualitatively unchanged across binary-vs-tri-state classifiers, and the headline VaR-uplift figures should be read as $K{=}2$-conditional only because the empirical $K{=}3$ ARI uses the binary highest-mean-state mapping by construction (\S{}2.2).

\subsection{Asymmetric per-asset persistence calibration}\label{app:asym_baseline}
The calibrated baseline above sweeps a \emph{symmetric} persistence pair ($P_{12}=P_{21}$ in the sensitivity rows). Under the post-revision pipeline, empirical 1h crisis shares span 16.3\% (USD/JPY) to 70.0\% (SPY); we therefore recalibrate the DGP per asset so the simulated stationary crisis share matches the empirical 1h share, while keeping the total transition rate at 5m fixed at the canonical ML-fit value $\tau \equiv P_{12}+P_{21} = 0.049$ (matrix-root conversion of the 1h fit; \S\ref{app:rss_sim}). For target stationary crisis share $\pi$ this gives $P_{12}=\pi\,\tau$ and $P_{21}=(1-\pi)\,\tau$. Two hundred canonical-pipeline replications per asset on the Gaussian-MS DGP yield:

% source: outputs/simulation_rss_asym_persistence_1h_anchor.csv + outputs/simulation_rss_asym_persistence_1d_anchor.csv
\begin{table}[H]
\centering
\small
\setlength{\tabcolsep}{4pt}
\caption{Asymmetric-persistence Gaussian-MS baseline calibrated to per-asset 1h crisis shares ($\tau=0.049$ ML-fit). Shares are computed at runtime from the post-revision pipeline (\texttt{outputs/simulation\_rss\_asym\_persistence\_1h\_anchor.csv}).}
\label{tab:asym_baseline}
\resizebox{\textwidth}{!}{%
\begin{tabular}{lcccccc}
\toprule
Asset & $\pi$ (1h) & $P_{12}$ & $P_{21}$ & All-4 mean ARI [IQR] & Empirical 2026 & Excess (emp $-$ cal) \\
\midrule
SPY     & 0.700 & 0.0343  & 0.0147  & 0.118\; [0.100, 0.131] & 0.183 & $+$0.07 \\
QQQ     & 0.619 & 0.0303  & 0.0187  & 0.117\; [0.098, 0.131] & 0.157 & $+$0.04 \\
USD/JPY & 0.163 & 0.0080  & 0.0410  & 0.134\; [0.110, 0.152] & 0.105 & $-$0.03 \\
CL      & 0.377 & 0.0185  & 0.0306  & 0.111\; [0.096, 0.125] & 0.394 & $+$0.28 \\
GLD     & 0.266 & 0.0131  & 0.0360  & 0.115\; [0.101, 0.128] & 0.155 & $+$0.04 \\
\bottomrule
\end{tabular}%
}
\par\medskip
{\footnotesize\justifying
\textit{Notes:} Per-asset DGP set so the stationary distribution of the 5m Markov chain matches the empirical 1h crisis share at $\tau=0.049$ (ML-fit, post-revision pipeline). 1h shares are computed at runtime by \texttt{compute\_panel\_crisis\_shares} (\texttt{src/experiments/exp\_16\_asym\_persistence\_baseline.py}) so the calibration target tracks the current pipeline output rather than a frozen snapshot. 200 canonical-pipeline replications per asset on 80-day 5m RTH paths; same Gaussian innovations as Suppl~\S\ref{app:rss_sim}.
\par}
\end{table}

\noindent The asymmetric baseline lies in 0.111--0.134 across all five assets, comparable to the symmetric calibrated reference (0.099--0.125). The cross-frequency ARI under Markov aggregation is governed primarily by regime-duration structure ($1/\tau$) and the GMM identifiability of two well-separated mass clusters, both of which are preserved under any redistribution of stationary mass between the two states; the per-asset asymmetric refinement does not move the baseline materially. The empirical-vs-calibrated gap is asset-specific: SPY/QQQ/CL/GLD have empirical ARI above their per-asset asymmetric baseline (excess $+0.04$ to $+0.28$, with CL the largest residual reflecting the 1d-touching-pair excess discussed in the main text), while USD/JPY's empirical 0.105 sits below its per-asset baseline 0.134 ($-0.03$ deviation). The empirical-coherence-beyond-DGP signal is therefore concentrated in SPY/QQQ/CL/GLD and absent for USD/JPY, consistent with the cross-asset inversion noted in the main text (\S{}3.2).

\paragraph{Per-asset full ML calibration baseline}
\label{app:per_asset_baseline}
The default calibrated baseline (Suppl~\S\ref{app:rss_sim}) ML-fits SPY 1h returns and uses that single $(P_{12}, P_{21}, \mu_0, \mu_1, \sigma_0, \sigma_1)$ tuple as the canonical reference for every empirical asset; the asymmetric-persistence baseline above pins $\tau$ at the SPY fit and only varies $\pi$. The third robustness check gives each asset a fully-independent ML fit on its own 1h log returns: per-asset $(P_{12}, P_{21})$ at the 1h scale, then $1/12$-th matrix root to 5m, then 200 canonical-pipeline replications per asset. Table~\ref{tab:per_asset_baseline} reports the result.

% source: outputs/simulation_rss_per_asset.csv (run via `python run.py extended_per_asset_baseline`)
\begin{table}[!htbp]
\centering
\small
\setlength{\tabcolsep}{4pt}
\caption{Per-asset full ML calibration baseline (2026 episode only; the 2022 OOS episode is excluded because QQQ and CL 2022 5m bars are absent from this revision round). Each asset's MS-Gaussian DGP is independently fit on its own 1h log returns; the resulting calibrated 5m DGP is run through the canonical pipeline $n{=}200$ times. Empirical ARI values from Table~1.}
\label{tab:per_asset_baseline}
\resizebox{\textwidth}{!}{%
\begin{tabular}{lccccc}
\toprule
Asset & $P_{12}^{1h}$ & $P_{21}^{1h}$ & $\tau^{1h}{=}P_{12}^{1h}{+}P_{21}^{1h}$ & Per-asset baseline mean [IQR] & Empirical (Gap) \\
\midrule
SPY     & 0.152 & 0.177 & 0.329 & 0.126\; [0.107, 0.139] & 0.183 ($+0.04$) \\
QQQ     & 0.179 & 0.227 & 0.406 & 0.117\; [0.101, 0.128] & 0.157 ($+0.03$) \\
USD/JPY & 0.051 & 0.343 & 0.394 & 0.108\; [0.078, 0.132] & 0.105 ($-0.03$) \\
CL      & 0.052 & 0.107 & 0.158 & 0.196\; [0.163, 0.221] & 0.394 ($+0.17$) \\
GLD     & 0.055 & 0.216 & 0.271 & 0.154\; [0.129, 0.171] & 0.155 ($-0.02$) \\
\bottomrule
\end{tabular}%
}
\par\medskip
{\footnotesize\justifying
\textit{Notes:} 1h-fit MS-Gaussian parameters from \texttt{statsmodels.MarkovRegression} with switching constant and variance, 50 EM iterations, 20 search restarts, fixed RNG seed (\texttt{src/core/calibration.py}). $\tau^{1h}$ ranges from 0.16 (CL: slow-switching commodity) to 0.41 (QQQ: fast-switching equity). Per-asset baselines are larger than the SPY-anchored baseline (0.113 IQR $[0.099, 0.125]$) wherever an asset has a slower switching rate than SPY. \emph{Three of five assets exceed their own per-asset IQR upper bound at the 4-frequency level} (SPY, QQQ, CL); USD/JPY and GLD lie within their per-asset IQR. All five assets exceed their per-asset \emph{intraday-only} IQR upper bound, preserving the headline ``uniform within-intraday excess'' result of \S{}3.1 under the strictest per-asset calibration.
\par}
\end{table}

\noindent The per-asset baseline answers ``is the SPY-1h calibrated reference generalisable to the rest of the panel?'' empirically. For session-driven ETFs (SPY, QQQ) and 24h FX (USD/JPY) the asset-specific baseline lies within $\pm 0.02$ of the SPY-1h reference (SPY 0.126, QQQ 0.117, USD/JPY 0.108 vs.\ SPY-1h 0.113). For commodity futures (CL) the per-asset baseline is materially higher (0.196) because CL's $\tau^{1h}=0.158$ is roughly half the SPY-fit $\tau^{1h}=0.329$, implying twice the regime persistence and correspondingly higher cross-frequency consistency under Markov aggregation. Even with this stricter per-asset reference, CL's empirical 0.394 still exceeds its per-asset IQR upper bound by $+0.17$, leaving the qualitative ``CL is the panel outlier'' finding intact under any calibration choice. GLD's 4-frequency empirical lies within its per-asset IQR; the gap to the SPY-anchored baseline (0.155 vs.\ 0.113, $+0.04$) was a $\tau$ mismatch, not a within-intraday excess (the intraday-only excess persists for GLD as well).

\paragraph{EM-restart placebo for the conservative-resolution VaR uplift}
\label{app:em_restart_placebo}
Because $\max(|\text{VaR}^{1h}|,|\text{VaR}^{1d}|) \ge |\text{VaR}^{1d}|$ by construction, a fraction of the headline 4.2\%--21.0\% range (Table~\ref{tab:var_uplift}) could in principle be the algebraic floor of taking the max of two correlated tail estimates rather than a resolution-specific signal. To bound that floor we apply the same $\max$-rule to two EM-restart variants of \emph{one} classifier on identical data: GMM is fit at one frequency with two random seeds (42 and 2024), regime-conditional 99\% empirical VaRs are computed under each variant separately, and the day-level $\max$-rule is applied across the two same-frequency classifiers. The resulting uplift carries no resolution information and bounds the algebraic-floor / classifier-noise component of the headline reading. Table~\ref{tab:em_restart_placebo} reports the placebo at both candidate probe frequencies (1h and 1d) under the canonical EM setting ($n_{\text{init}}=10$, identical to the main pipeline) and an adversarial classifier-noise stress test ($n_{\text{init}}=1$, where each EM run uses a single random initialisation and is therefore maximally sensitive to seed).

% source: outputs/var_uplift_em_restart_placebo.csv          (2026)
%         outputs_2022/var_uplift_em_restart_placebo.csv     (2022 OOS)
\begin{table}[H]
\centering
\small
\setlength{\tabcolsep}{4pt}
\caption{EM-restart placebo for the conservative-resolution VaR uplift: $\max(|\text{VaR}_A|,|\text{VaR}_B|)$ where $A$ and $B$ are two-seed EM-restart variants of one classifier at the same probe frequency. Headline 1h-vs-1d uplift (Table~\ref{tab:var_uplift}) repeated for direct comparison.}
\label{tab:em_restart_placebo}
\resizebox{\textwidth}{!}{%
\begin{tabular}{llcccccc}
\toprule
 & & \multicolumn{2}{c}{Placebo $n_{\text{init}}=10$ (canonical)} & \multicolumn{2}{c}{Placebo $n_{\text{init}}=1$ (stress)} & \multicolumn{2}{c}{Headline 1h-vs-1d} \\
\cmidrule(lr){3-4}\cmidrule(lr){5-6}\cmidrule(lr){7-8}
Episode & Asset & 1h uplift & 1d uplift & 1h uplift & 1d uplift & Uplift & Disagree (\%) \\
\midrule
2026 & SPY     & 0.00\% & 0.00\% & 0.00\% & 0.00\% & 7.5\%  & 26.2 \\
2026 & QQQ     & 0.00\% & 0.00\% & 1.20\% & 4.60\% & 4.2\%  & 27.0 \\
2026 & USD/JPY & 0.00\% & 0.00\% & 0.00\% & 6.55\% & 7.8\%  & 52.6 \\
2026 & CL      & 0.00\% & 0.00\% & 0.00\% & 0.00\% & 21.0\% & 17.6 \\
2026 & GLD     & 0.00\% & 0.00\% & 0.00\% & 0.00\% & 13.9\% & 63.1 \\
2022 & SPY     & 0.00\% & 0.00\% & 0.00\% & 3.97\% & 10.0\% & 43.9 \\
2022 & USD/JPY & 0.15\% & 0.00\% & 0.33\% & 0.00\% & 12.3\% & 41.6 \\
2022 & GLD     & 0.00\% & 0.00\% & 0.49\% & 0.00\% & 12.7\% & 49.6 \\
\bottomrule
\end{tabular}%
}
\par\medskip
{\footnotesize\justifying
\textit{Notes:} Two EM-restart variants of the GMM are fit at the placebo frequency with seeds 42 and 2024; both use the canonical \texttt{sklearn.GaussianMixture} estimator (full covariance, k-means++ init) on the same log-volatility feature matrix. Placebo uplift is the full-sample mean of $(\max(|\text{VaR}_A|,|\text{VaR}_B|)/|\text{VaR}_A|) - 1$, computed identically to the headline column except both VaRs come from same-frequency restart variants rather than 1h / 1d classifiers. \emph{Canonical $n_{\text{init}}=10$ matches the main-pipeline setting} (Section~2.2); under that classifier the two seeds produce identical full-sample VaR uplifts on every (asset, frequency) cell except 2022 USD/JPY 1h, where they differ on 0.2\% of 5m bars and produce a 0.15\% full-sample uplift. The non-trivial cells under the adversarial $n_{\text{init}}=1$ stress test are 2026 USD/JPY 1d (6.55\%) and 2026 QQQ 1d (4.60\%), reflecting medium 1d component overlaps (0.582 / 0.559, Table~\ref{tab:gmm_diag}) and the smallest sample size at 1d ($n=124$--$156$). All other cells under the stress test return $\le 1.20\%$ uplift. \emph{Implication.} Under the classifier the paper actually uses, the algebraic-floor component of the headline 4.2\%--21.0\% range (2026) and 10.0\%--12.7\% range (2022 OOS) is at most 0.15\% on any cell; the headline range is therefore essentially fully attributable to cross-resolution dissonance rather than to the algebraic floor of the $\max$-rule.
\par}
\end{table}

\section{Resolution-dissonance heatmaps for QQQ, CL, USD/JPY, and GLD}\label{app:timeline_suppl}

Figures~\ref{fig:timeline_suppl_a} and~\ref{fig:timeline_suppl_b} reproduce the resolution-dissonance heatmap of \S{}3.3 for the four remaining 2026 assets, split across two figures so each pair fits cleanly on one page. Reading conventions are identical to Figure~1; the layout caveat differs by asset.

\begin{figure}[!htbp]
\centering
\IfFileExists{QQQ_timeline.png}{%
  \includegraphics[width=0.95\linewidth]{QQQ_timeline.png}\par\smallskip
}{%
  \fbox{\parbox{0.85\linewidth}{\centering\small [QQQ\_timeline.png]}}\par\smallskip
}
\IfFileExists{CL_timeline.png}{%
  \includegraphics[width=0.95\linewidth]{CL_timeline.png}%
}{%
  \fbox{\parbox{0.85\linewidth}{\centering\small [CL\_timeline.png]}}%
}
\caption{\textbf{Resolution-dissonance heatmaps for QQQ and CL (2025-11 to 2026-05).} Same reading conventions as Figure~1. \textbf{QQQ (top):} all four rows are GMM-derived (none degenerate); the 5m--15m agreement is high while the 15m$\to$1h drop mirrors SPY, consistent with the equity-ETF pattern (\S{}3.1). \textbf{CL (bottom):} the 5m and 15m rows reflect percentile-fallback labels (top 20\% of log-volatility = crisis) because the underlying GMM is degenerate at those resolutions (\S{}2.2, Table~\ref{tab:gmm_diag}); the load-bearing within-CL evidence is the 1h vs.\ 1d row contrast, where both fits are well-identified and exhibit a $\sim$4-month \emph{episode-onset gap} (1d row red from Nov 2025 onwards tracking broadly elevated daily volatility; 1h row blue until late February 2026, then red as supply-shock-spike weeks begin), with the same days inside the spike season classified as crisis at 1h being absorbed into 1d-calm bars on individual single-day spikes such as 6 March 2026 (\S{}3.2).}
\label{fig:timeline_suppl_a}
\end{figure}

\begin{figure}[!htbp]
\centering
\IfFileExists{USDJPY_timeline.png}{%
  \includegraphics[width=0.95\linewidth]{USDJPY_timeline.png}\par\smallskip
}{%
  \fbox{\parbox{0.85\linewidth}{\centering\small [USDJPY\_timeline.png]}}\par\smallskip
}
\IfFileExists{GLD_timeline.png}{%
  \includegraphics[width=0.95\linewidth]{GLD_timeline.png}%
}{%
  \fbox{\parbox{0.85\linewidth}{\centering\small [GLD\_timeline.png]}}%
}
\caption{\textbf{Resolution-dissonance heatmaps for USD/JPY and GLD (2025-11 to 2026-05).} Same reading conventions as Figure~1. \textbf{USD/JPY (top):} the 15m row reflects percentile-fallback labels (degenerate GMM); the 5m row's GMM is non-degenerate but at the unimodality boundary (overlap 0.85, regime-concept failure, \S{}3.1). \textbf{GLD (bottom):} all four rows are GMM-derived (none degenerate), but 5m and 15m have high overlap (0.59, 0.71) reflecting weak vol-regime separation at fine resolutions. In both assets the 1d row decouples sharply from the intraday rows on a large fraction of days, consistent with the 18--63\% 1h--1d disagreement rates in Table~\ref{tab:var_uplift}.}
\label{fig:timeline_suppl_b}
\end{figure}

\clearpage  % flush all pending floats so References do not appear before figures
\bibliographystyle{apalike}
\bibliography{refs}

\end{document}
